\hebrewsection{Case Studies and Practical Implementations}

\subsection{Case Study: The GSM Standard}
GSM (Global System for Mobile Communications) was the first widely adopted digital cellular standard. It uses a combination of TDMA (Time Division Multiple Access) and FDMA (Frequency Division Multiple Access).

\subsubsection{Modulation in GSM: GMSK}
As discussed in Chapter 3, GSM uses Gaussian Minimum Shift Keying. GMSK provides a constant envelope, which is critical for efficient power amplifier operation in mobile handsets. The Gaussian filter bandwidth-time product ($BT$) is typically set to 0.3.

\subsection{Case Study: Digital Subscriber Line (DSL)}
DSL technology provides high-speed internet over traditional copper telephone lines. It uses \textbf{Discrete Multi-Tone (DMT)} modulation, which is a form of OFDM.

\subsubsection{Bit Loading}
DMT adaptively allocates bits to different subcarriers based on the Signal-to-Noise Ratio (SNR) of each subchannel. This allows DSL to maximize the data rate even in the presence of frequency-selective impairments like bridge taps.

\subsection{Case Study: 5G New Radio (NR)}
5G NR introduces several advanced features to support massive machine-type communications (mMTC) and ultra-reliable low-latency communications (URLLC).

\subsubsection{Massive MIMO and Beamforming}
5G uses large antenna arrays (up to 256 elements) to perform highly directional beamforming. This increases the signal strength at the intended user while reducing interference for others.

\subsection{Case Study: Evolution of IEEE 802.11 (Wi-Fi)}
The Wi-Fi standard has evolved from 2 Mbps (802.11b) to multi-Gbps rates (802.11ax/be) by increasing bandwidth, using higher-order modulation (1024-QAM), and employing massive MIMO.

\subsubsection{Wi-Fi 6 (802.11ax)}
Wi-Fi 6 introduced \textbf{OFDMA} (Orthogonal Frequency Division Multiple Access), allowing the router to serve multiple clients simultaneously within a single symbol interval by allocating subcarrier subsets (Resource Units) to different users.

\subsection{Practical Implementation: FPGA-based Digital Down-Converter (DDC)}
A DDC is used in modern receivers to shift a high-frequency IF signal to baseband and reduce the sampling rate.
\begin{enumerate}
    \item \textbf{Numerical Control Oscillator (NCO):} Generates a complex sinusoid.
    \item \textbf{Mixer:} Performs the frequency shift.
    \item \textbf{CIC Filter:} Efficiently performs decimation and initial filtering.
    \item \textbf{FIR Decimation Filter:} Provides sharp filtering and further decimation.
\end{enumerate}

\begin{pythonbox}{Simple Digital Down-Conversion}
\begin{verbatim}
import numpy as np
from scipy import signal

# IF Signal at 10 MHz, Sampled at 100 MHz
fs = 100e6
t = np.arange(1000) / fs
x_if = np.cos(2*np.pi*10e6*t)

# NCO at 10 MHz
nco = np.exp(-1j * 2*np.pi*10e6*t)

# Mixing to Baseband
x_baseband = x_if * nco

# Low-pass filter and Decimate by 10
b = signal.firwin(64, 1e6, fs=fs)
x_lowpass = signal.lfilter(b, 1, x_baseband)
x_final = x_lowpass[::10]
\end{verbatim}
\end{pythonbox}
